% GENERATED FILE. Edit canonical lessons, then run npm run generate:books.
% canonical-source-hash: fnv1a64:451fc8d56b57639a
% canonical-lessons: SA-C159-jirakam, SA-C159-sarsapah, SA-C159-hingu, SA-C159-tilah, SA-C159-canakah

\chapter{Cumin, Mustard, Asafoetida, Sesame, Chickpea}
\label{ch:sa-cumin-mustard-asafoetida-sesame-chickpea}
\hlchaptermodality{159}

\begin{chapteropening}
\textbf{By the end of this chapter:} \emph{I can say cumin, mustard, asafoetida, sesame and a chickpea.}

Complete the last lesson of chapter 159: I can say cumin, mustard, asafoetida, sesame and a chickpea.
\end{chapteropening}

\section[\texorpdfstring{\sk{जीरकम्} (jīrakam)}{जीरकम् (jīrakam)}]{\sk{जीरकम्} (jīrakam) — cumin}

\label{lesson:SA-C159-jirakam}



\begin{quote}
Before the new one: say the Sanskrit for a small lamp, then the Sanskrit for a door panel, a shutter.
\end{quote}

\subsection*{You'll want to know: \sk{जीरकम्}}
\textbf{\sk{जीरकम्}} — \emph{jīrakam} — \textquotedblleft{}cumin\textquotedblright{}.

\begin{grammarlens}[title={What you've built}]
One more word for this chapter.
\end{grammarlens}

\subsection*{Your turn}
\begin{itemize}
\raggedright
  \item \emph{Say it:} \emph{jīrakam}
  \item \emph{Say it:} \emph{jīrakam}, once more
  \item \emph{Say it:} say \emph{jīrakam}
  \item \emph{Recall:} say the Sanskrit for a small lamp, then the Sanskrit for a door panel, a shutter, then say \emph{jīrakam} again
\end{itemize}

\begin{tcolorbox}[breakable,colback=teal!4,colframe=teal!35!black,title={Before you move on}]
What does \sk{जीरकम्} mean? (\textquotedblleft{}Cumin\textquotedblright{}.) Say it once more.
\end{tcolorbox}
\section[\texorpdfstring{\sk{सर्षपः} (sarṣapaḥ)}{सर्षपः (sarṣapaḥ)}]{\sk{सर्षपः} (sarṣapaḥ) — mustard}

\label{lesson:SA-C159-sarsapah}



\begin{quote}
Before the new one: say the Sanskrit for a door panel, a shutter, then the Sanskrit for cumin.
\end{quote}

\subsection*{You'll want to know: \sk{सर्षपः}}
\textbf{\sk{सर्षपः}} — \emph{sarṣapaḥ} — \textquotedblleft{}mustard\textquotedblright{}.

\begin{grammarlens}[title={What you've built}]
One more word for this chapter.
\end{grammarlens}

\subsection*{Your turn}
\begin{itemize}
\raggedright
  \item \emph{Say it:} \emph{sarṣapaḥ}
  \item \emph{Say it:} \emph{sarṣapaḥ}, once more
  \item \emph{Say it:} say \emph{sarṣapaḥ}
  \item \emph{Recall:} say the Sanskrit for a door panel, a shutter, then the Sanskrit for cumin, then say \emph{sarṣapaḥ} again
\end{itemize}

\begin{tcolorbox}[breakable,colback=teal!4,colframe=teal!35!black,title={Before you move on}]
What does \sk{सर्षपः} mean? (\textquotedblleft{}Mustard\textquotedblright{}.) Say it once more.
\end{tcolorbox}
\section[\texorpdfstring{\sk{हिंगु} (hiṅgu)}{हिंगु (hiṅgu)}]{\sk{हिंगु} (hiṅgu) — asafoetida}

\label{lesson:SA-C159-hingu}



\begin{quote}
Before the new one: say the Sanskrit for cumin, then the Sanskrit for mustard.
\end{quote}

\subsection*{You'll want to know: \sk{हिंगु}}
\textbf{\sk{हिंगु}} — \emph{hiṅgu} — \textquotedblleft{}asafoetida\textquotedblright{}.

\begin{grammarlens}[title={What you've built}]
One more word for this chapter.
\end{grammarlens}

\subsection*{Your turn}
\begin{itemize}
\raggedright
  \item \emph{Say it:} \emph{hiṅgu}
  \item \emph{Say it:} \emph{hiṅgu}, once more
  \item \emph{Say it:} say \emph{hiṅgu}
  \item \emph{Recall:} say the Sanskrit for cumin, then the Sanskrit for mustard, then say \emph{hiṅgu} again
\end{itemize}

\begin{tcolorbox}[breakable,colback=teal!4,colframe=teal!35!black,title={Before you move on}]
What does \sk{हिंगु} mean? (\textquotedblleft{}Asafoetida\textquotedblright{}.) Say it once more.
\end{tcolorbox}
\section[\texorpdfstring{\sk{तिलः} (tilaḥ)}{तिलः (tilaḥ)}]{\sk{तिलः} (tilaḥ) — sesame}

\label{lesson:SA-C159-tilah}



\begin{quote}
Before the new one: say the Sanskrit for mustard, then the Sanskrit for asafoetida.
\end{quote}

\subsection*{You'll want to know: \sk{तिलः}}
\textbf{\sk{तिलः}} — \emph{tilaḥ} — \textquotedblleft{}sesame\textquotedblright{}.

\begin{grammarlens}[title={What you've built}]
One more word for this chapter.
\end{grammarlens}

\subsection*{Your turn}
\begin{itemize}
\raggedright
  \item \emph{Say it:} \emph{tilaḥ}
  \item \emph{Say it:} \emph{tilaḥ}, once more
  \item \emph{Say it:} say \emph{tilaḥ}
  \item \emph{Recall:} say the Sanskrit for mustard, then the Sanskrit for asafoetida, then say \emph{tilaḥ} again
\end{itemize}

\begin{tcolorbox}[breakable,colback=teal!4,colframe=teal!35!black,title={Before you move on}]
What does \sk{तिलः} mean? (\textquotedblleft{}Sesame\textquotedblright{}.) Say it once more.
\end{tcolorbox}
\section[\texorpdfstring{\sk{चणकः} (caṇakaḥ)}{चणकः (caṇakaḥ)}]{\sk{चणकः} (caṇakaḥ) — a chickpea}

\label{lesson:SA-C159-canakah}



\begin{quote}
Before the new one: say the Sanskrit for asafoetida, then the Sanskrit for sesame.
\end{quote}

\subsection*{You'll want to know: \sk{चणकः}}
\textbf{\sk{चणकः}} — \emph{caṇakaḥ} — \textquotedblleft{}a chickpea\textquotedblright{}.

\begin{grammarlens}[title={What you've built}]
One more word for this chapter.
\end{grammarlens}

\subsection*{Your turn}
\begin{itemize}
\raggedright
  \item \emph{Say it:} \emph{caṇakaḥ}
  \item \emph{Say it:} \emph{caṇakaḥ}, once more
  \item \emph{Say it:} say \emph{caṇakaḥ}
  \item \emph{Recall:} say the Sanskrit for asafoetida, then the Sanskrit for sesame, then say \emph{caṇakaḥ} again
\end{itemize}

\begin{tcolorbox}[breakable,colback=teal!4,colframe=teal!35!black,title={Before you move on}]
What does \sk{चणकः} mean? (\textquotedblleft{}A chickpea\textquotedblright{}.) Say it once more.
\end{tcolorbox}
